\section{Why an agent is not a rank}
\label{sec:mismatch}

The interesting part of this work is not the parts of MPI that transfer
unchanged---those are numerous but unsurprising---but the four places where
the executor's nature forces the design to differ. We state each mismatch
precisely, because each one is the justification for a specific protocol
feature in \Cref{sec:design}, and because a design that borrows MPI's
vocabulary without acknowledging them would inherit the wrong semantics
along with the right ones.

\subsection{M1: context is consumed, not occupied}
\label{sec:m1}

An MPI receive buffer is reusable. A rank can receive a terabyte over a
job's lifetime while never holding more than a megabyte at once, because
the memory occupied by message $k$ is available again for message $k+1$.
Memory pressure is a function of the \emph{peak live set}, and the standard
consequence is that a collective's memory cost is analysed per round.

Every token an agent reads enters its context and stays there. A rank that
receives 200 messages of 5{,}000 tokens has spent a million tokens whether
or not it needed two of them simultaneously. Pressure is a function of
\emph{cumulative ingest}. Three consequences follow, and all three are
protocol-visible:

\begin{itemize}
\item A receive is not a data movement but an irreversible allocation
against a quota, so it needs \emph{admission control}: the runtime must
decide whether a payload may enter before materialising it
(\Cref{sec:design-p2p}).
\item Continuing past the quota requires rewriting history---replacing a
span of consumed context with a smaller summary of it. This is a garbage
collector whose liveness analysis is semantic, so the protocol supplies the
mechanism and accounting and the harness supplies the policy.
\item Collective feasibility becomes a first-class question. An MPI
collective is always feasible; it may be slow, but no rank is \emph{unable}
to hold its share. \op{Allgather} over $p$ agents requires every rank to
ingest $(p-1)s$ tokens by definition, so it stops working at
$p > C/s$ regardless of algorithm. The runtime must be able to say so before
the job spends anything.
\end{itemize}

There is a second-order effect that surprised us and that
\Cref{sec:cost-cumulative} develops: because ingest accumulates across
\emph{rounds}, the depth of a reduction tree is not free. An MPI rank that
receives $k-1$ buffers per round reuses the same memory every round, so its
peak is one round's worth and depth costs only latency. An agent root of a
depth-$d$ tree pays $(k-1)d$ contributions over the reduction's life.
Planning a tree against the per-round figure produces schedules that look
feasible, run two rounds, and then exhaust the root.

\subsection{M2: reduction operators are semantic}
\label{sec:m2}

\mpiop{SUM} is associative, commutative, and output-size-preserving.
Floating-point addition is not exactly associative, and MPI's community has
long accepted the resulting last-bit non-reproducibility as the price of
letting the implementation reassociate freely.

Agent operators fail all three properties, and the failures are not in the
last bit:

\begin{itemize}
\item ``Concatenate in rank order'' is associative but not commutative.
\item ``Take a majority vote'' is commutative but \emph{not} associative: a
pairwise tree of votes computes something different from a flat vote over
the same inputs, and which answer wins depends on the number of ranks.
\item ``Summarise these two summaries'' is approximately associative at
best, and its output must be strictly smaller than its input or a tree of
depth $d$ produces an output of size $2^d$.
\item ``Merge these two drafts'' is none of the above.
\end{itemize}

So a protocol that wants to choose collective schedules on the
application's behalf must be \emph{told} the operator's algebra, and must
refuse schedules the algebra does not license. \mpiop{Op\_create} already
takes a \code{commute} flag for exactly this reason; \ampi{} needs that flag
and two more---idempotency, which determines tolerance of duplicate
delivery, and an output bound, which determines the legal tree depth
(\Cref{sec:ops}).

A related consequence: operator application is itself an agent turn. It
costs seconds and money and can fail. Reductions therefore need retry,
memoisation, and contract validation at every internal node---machinery that
in MPI would be absurd and here is unavoidable.

\subsection{M3: failure is not fail-stop}
\label{sec:m3}

ULFM's model is binary: a rank is alive or it has failed, and failure is
detected because the process is gone. Agents fail in at least four
distinguishable ways, and the distinctions matter because the correct
response differs:

\begin{description}[leftmargin=1.4em,style=nextline]
\item[Crashed] the process is gone. Detected by absence of heartbeat;
respond by replacing or shrinking.
\item[Stalled] the process is healthy and making no progress---looping on a
tool call, retrying a rate-limited endpoint, or re-reading its own output.
Indistinguishable from a working agent by any liveness test, and
indistinguishable from a \emph{slow} agent without a progress signal.
\item[Drifted] the process returns output that violates the contract it was
given: wrong schema, wrong language, a summary where a patch was expected.
Perfectly alive, perfectly prompt, and useless. Detected only by validating
content.
\item[Bankrupt] the process has exhausted its token or currency budget.
Deterministic, predictable, and invisible to every failure detector in the
distributed-systems literature, because no classical process runs out of
money.
\end{description}

Conflating these produces the two pathologies every practitioner reports:
killing a slow-but-working agent because it looked dead, and waiting forever
for a fast-but-broken one because it looked alive. The protocol therefore
needs two independent timeouts---liveness and progress---and a
content-level error class alongside the transport-level ones.

There is also an asymmetry in the other direction, and it is favourable.
Restarting an MPI rank is expensive: its state is the application's memory,
which is why coordinated checkpointing and the entire literature on
message-logging and the domino effect
exist~\cite{Elnozahy2002,Chandy1985}. An agent rank's recoverable state is
its role, its checkpoint, and the messages it has acknowledged---kilobytes.
So the balance between ULFM's \emph{shrink} (cheap for the runtime,
expensive for the application, whose data distribution was written against
the old indices) and FT-MPI's \emph{rebuild} (expensive for the runtime,
free for the application) tips the other way, and \Cref{sec:ft} makes
replacement the recommended path.

Finally, the classical recovery technique of message-log replay depends on
the piecewise-deterministic assumption: re-executing a process from a
checkpoint with the same message sequence reproduces its
behaviour~\cite{Alvisi1995}. Agents violate this outright. Our substitute is
content-addressed memoisation, which gives up bitwise reproduction and keeps
the property that actually matters: a re-executed rank tells its peers the
same thing it told them before.

\subsection{M4: the cost model inverts}
\label{sec:m4}

In HPC, $\alpha \approx 1\,\mu s$, $\beta^{-1} \approx 10\,\mathrm{GB/s}$,
and a floating-point operation costs a fraction of a nanosecond. Every
collective algorithm in \Cref{sec:background} is a trade between the
$\alpha$ term and the $\beta$ term, and computation is nowhere in the
analysis because it is free relative to communication.

For agents: transmitting a message is a file write or an HTTP round
trip---microseconds to milliseconds. Producing the message is an inference
call: tens of seconds and, at current prices, cents to dollars. The ratio
between the cost of moving a payload and the cost of creating it is between
$10^5$ and $10^7$, and it points the opposite way.

Three consequences:

\begin{enumerate}
\item \textbf{Latency-optimal algorithms win essentially everywhere.} The
HPC crossover at which bandwidth-optimal ring algorithms overtake
logarithmic ones does not appear, because the bandwidth term is negligible.
A ring allreduce over 64 agents costs 126 sequential turns; recursive
doubling costs 6. \Cref{sec:eval-micro} measures this.
\item \textbf{The right unit of cost is the turn, not the second.} Wall
clock varies by model, provider and time of day; rounds on the critical path
do not. We report both, and treat rounds as the invariant.
\item \textbf{Redundant computation is the thing to avoid.} Chatty protocols
are fine; recomputing an agent's output is not. This is why memoisation is a
protocol concern here and a curiosity in MPI.
\end{enumerate}

\subsection{What the mismatches do not change}

It is worth being explicit about the negative result, because it is the
paper's main structural claim. None of M1--M4 touches the matching
semantics, the communicator/context construction, the group algebra, the
topology interface, the collective \emph{catalogue}, the epoch structure of
one-sided operations, the two-phase I/O idea, the shape of
revoke/shrink/agree, or the profiling interface. Those transfer verbatim.
The mismatches bear on the type system, the operator algebra, algorithm
selection, and the failure lattice---four well-localised places. A
message-passing interface for agents is mostly MPI.
